\documentclass[11pt]{amsart}

\usepackage[T1]{fontenc}
\usepackage{lmodern}
\usepackage{microtype}
\usepackage{mathtools,amssymb,amsthm}
\usepackage[hidelinks]{hyperref}

\hypersetup{
  pdftitle={A Kirkman Triple System of Order 27 with a Steiner Triple System of
            Order 7 as a Subdesign}
}

\newtheorem{theorem}{Theorem}
\newtheorem{lemma}[theorem]{Lemma}
\newtheorem{proposition}[theorem]{Proposition}
\newtheorem{corollary}[theorem]{Corollary}
\theoremstyle{remark}
\newtheorem*{remark}{Remark}

\newcommand{\KTS}{\mathrm{KTS}}
\newcommand{\STS}{\mathrm{STS}}

\title[A $\KTS(27)$ with an $\STS(7)$ subdesign]
{A Kirkman Triple System of Order $27$ with a Steiner Triple System of
Order $7$ as a Subdesign}

\date{}

\subjclass[2020]{05B07, 05B05, 05B30}
\keywords{Kirkman triple system, Steiner triple system, subdesign,
resolvable design, Fano plane, automorphism of order three}

\begin{document}

\begin{abstract}
Dukes and Lamken write that, although a $\KTS(27)$ with a sub-$\STS(u)$ is easy
to construct for $u=3,9,13$, ``it appears difficult to find a $\KTS(27)$ with an
$\STS(7)$ subdesign''. We exhibit one. Its thirteen parallel classes are printed
in full, and everything claimed about them is elementary counting on the printed
table. Whether such a design has been exhibited before was not determined here;
no search of the literature or of design databases was carried out. The orders
$v=33$ and $v=39$ of the same sentence of Dukes and Lamken are untouched.
\end{abstract}

\maketitle

\section{The problem}

A \emph{Steiner triple system} $\STS(v)$ is a set of $v$ points together with a
collection of $3$-subsets, called blocks or triples, covering every pair of
points exactly once. It is a \emph{Kirkman triple system} $\KTS(v)$ if its
blocks can be partitioned into \emph{parallel classes}, each a partition of the
point set; there are then necessarily $(v-1)/2$ of them. A \emph{subdesign} of
order $u$ is a $u$-subset $W$ of the points together with the blocks lying
inside $W$, this being required to be an $\STS(u)$.

Dukes and Lamken \cite{DL} write:
\begin{quote}
There remain a number of interesting and difficult computational problems in
finding KTS with STS subdesigns. We can construct KTS($v$) with all possible
admissible STS($u$) for the two smallest cases, $v=15$ and $v=21$. But, although
it is easy to construct KTS($27$) with sub-STS($u$) for $u=3,9,13$, it appears
difficult to find a KTS($27$) with an STS($7$) subdesign. Similarly, the
existence of KTS($v$) with an STS($7$) subdesign is unknown for $v=33$ and $39$.
New direct constructions are needed for these cases.
\end{quote}
The passage is quoted from the e-print arXiv:2110.07874v1 of \cite{DL}.

Two conventions of \cite{DL} decide what counts as a witness. First, the
subdesign notion is the one recalled above: a point subset together with the
blocks inside it. Second, no compatibility with the resolution is imposed ---
\cite{DL} writes that ``a `Kirkman subsystem' or sub-KTS($u$) is a
subdesign which is resolvable in such a way that parallel classes of the KTS($u$)
are restrictions of parallel classes of the KTS($v$). We do not assume such
extra structure for subdesigns here.'' At $u=7$ the stronger notion cannot apply
in any case, since $3\nmid 7$ and a set of $7$ points has no partition into
triples.

The problem is attributed in \cite{DL} to Open Problem~4 of Stinson's survey
\cite{Stinson}, with \cite[\S19.7]{ColbournRosa} named there as the section
recording its status.

\section{The design}\label{sec:design}

Let $V=\{0,1,\ldots,26\}$ and $U=\{0,1,2,3,4,5,6\}$. Table~\ref{tab:kts}
lists thirteen sets $C_{00},\ldots,C_{12}$ of nine triples each, each set
written on two lines, the first ending in a vertical bar to mark the
continuation. Call the resulting collection of $13\cdot 9=117$ triples $\mathcal D$.

\begin{table}[htbp]
\caption{The thirteen parallel classes. Points are $0,\ldots,26$; triples within
a class are separated by \texttt{|}, and a trailing \texttt{|} means the class
continues on the next line.}
\label{tab:kts}
\small
\begin{verbatim}
C00: 0 1 3 | 2 17 18 | 4 11 21 | 5 22 23 | 6 12 19 |
     7 9 14 | 8 15 20 | 10 13 24 | 16 25 26
C01: 0 7 8 | 1 2 4 | 3 20 24 | 5 19 26 | 6 18 25 |
     9 10 11 | 12 15 22 | 13 16 23 | 14 17 21
C02: 0 11 16 | 1 19 25 | 2 3 5 | 4 9 20 | 6 7 24 |
     8 21 26 | 10 14 22 | 12 13 18 | 15 17 23
C03: 0 9 17 | 1 10 18 | 2 20 26 | 3 4 6 | 5 7 25 |
     8 22 24 | 11 12 23 | 13 14 19 | 15 16 21
C04: 0 4 5 | 1 16 20 | 2 10 23 | 3 14 18 | 6 21 22 |
     7 11 13 | 8 17 19 | 9 12 26 | 15 24 25
C05: 0 10 15 | 1 5 6 | 2 11 19 | 3 7 26 | 4 18 24 |
     8 23 25 | 9 13 21 | 12 14 20 | 16 17 22
C06: 0 2 6 | 1 9 22 | 3 21 23 | 4 15 19 | 5 13 20 |
     7 10 12 | 8 16 18 | 11 14 25 | 17 24 26
C07: 0 19 23 | 1 15 26 | 2 7 22 | 3 13 17 | 4 8 12 |
     5 14 24 | 6 9 16 | 10 21 25 | 11 18 20
C08: 0 20 21 | 1 8 13 | 2 16 24 | 3 12 25 | 4 7 23 |
     5 10 17 | 6 14 15 | 9 18 19 | 11 22 26
C09: 0 18 22 | 1 7 21 | 2 8 14 | 3 11 15 | 4 17 25 |
     5 12 16 | 6 13 26 | 9 23 24 | 10 19 20
C10: 0 14 26 | 1 12 17 | 2 9 25 | 3 10 16 | 4 13 22 |
     5 8 11 | 6 20 23 | 7 15 18 | 19 21 24
C11: 0 12 24 | 1 14 23 | 2 13 15 | 3 8 9 | 4 10 26 |
     5 18 21 | 6 11 17 | 7 16 19 | 20 22 25
C12: 0 13 25 | 1 11 24 | 2 12 21 | 3 19 22 | 4 14 16 |
     5 9 15 | 6 8 10 | 7 17 20 | 18 23 26
\end{verbatim}
\end{table}

\begin{theorem}\label{thm:main}
$\mathcal D$ is a Kirkman triple system of order $27$ with
$C_{00},\ldots,C_{12}$ as its parallel classes, and $U$ carries an $\STS(7)$
subdesign of $\mathcal D$. In particular a $\KTS(27)$ with an $\STS(7)$
subdesign exists.
\end{theorem}

\begin{proof}
Each of the thirteen lists in Table~\ref{tab:kts} contains nine triples whose
$27$ entries are $0,1,\ldots,26$ without repetition, so each $C_{ij}$ is a
partition of $V$ and $\mathcal D$ is the disjoint union of thirteen parallel
classes. The $117$ triples carry $117\cdot 3=351=\binom{27}{2}$ pairs, and the
pairs so obtained are all distinct; being $\binom{27}{2}$ in number they are
therefore all of them, each once. Hence $\mathcal D$ is an $\STS(27)$ and
$C_{00},\ldots,C_{12}$ is a resolution of it, i.e.\ a $\KTS(27)$; the number of
classes is $(27-1)/2=13$.

The triples of $\mathcal D$ contained in $U$ are exactly
\[
 \{0,1,3\},\ \{1,2,4\},\ \{2,3,5\},\ \{3,4,6\},\ \{0,4,5\},\ \{1,5,6\},\
 \{0,2,6\},
\]
one in each of $C_{00},C_{01},\ldots,C_{06}$ in that order. They carry
$7\cdot 3=21=\binom{7}{2}$ pairs, all distinct, hence all pairs inside $U$, each
once. So $U$ together with these seven triples is an $\STS(7)$: it is the Fano
plane, being the unique $\STS(7)$. By the definition recalled in \S1 it is a
subdesign of $\mathcal D$ of order $7$.
\end{proof}

\begin{remark}
No triple of $\mathcal D$ meets $U$ in exactly two points: such a pair lies
inside $U$ and is already covered by one of the seven triples above, and pairs
are covered once. So the $117$ triples split as
$7+70+40$: seven inside $U$, seventy meeting $U$ in one point, and forty disjoint from $U$. Concretely, each of
$C_{00},\ldots,C_{06}$ consists of one inner, four one-point and four disjoint
triples, and each of $C_{07},\ldots,C_{12}$ of seven one-point and two
disjoint triples.
\end{remark}

\section{The subdesigns of the printed design}

Write $\langle S\rangle$ for the \emph{closure} of a set $S$ of points: the
smallest subset containing $S$ and closed under the operation sending a pair of
its points to the third point of the block through them. A closed set of size at
least $3$, with the blocks inside it, is a subdesign.

\begin{lemma}\label{lem:2w1}
If $W$ is a subdesign of order $w\ge 3$ of an $\STS(v)$ and $W$ is proper, then
$v\ge 2w+1$.
\end{lemma}

\begin{proof}
Fix $x\notin W$. For $p\in W$ let $f(p)$ be the third point of the block through
$x$ and $p$. Then $f(p)\notin W$, since a block meeting $W$ in two points has
its third point in $W$ and would give a pair of $W$ covered twice; and
$f(p)\ne x$. If $f(p)=f(q)$ with $p\ne q$ then the block $\{x,p,f(p)\}$ contains
$q$, so $p,q,f(p)$ are collinear and $f(p)\in W$, a contradiction. Hence $f$ is
injective into $V\setminus(W\cup\{x\})$, and $v\ge w+1+w$.
\end{proof}

\begin{proposition}\label{prop:inv}
$\mathcal D$ has exactly one subdesign of order $7$, namely $U$, and none of
order $9$ or $13$. Consequently the number of order-$7$ subdesigns through a
point equals $1$ for each of the seven points of $U$ and $0$ for each of the
other twenty; hence $\operatorname{Aut}(\mathcal D)$ is not transitive on points.
\end{proposition}

\begin{proof}
Let $W$ be a subdesign of $\mathcal D$ of order $u\in\{7,9,13\}$ and let
$x,y,z\in W$ be three points not on a common block; such a triple exists because
$u\ge 7$. Then $\langle x,y,z\rangle\subseteq W$ is a subdesign of some order
$u'$ with $3<u'\le u$, and $u'$ is an admissible order for a Steiner triple
system, hence $u'\in\{7,9,13\}$ with $u'\le u$. If $u'<u$ then
Lemma~\ref{lem:2w1} forces $u\ge 2u'+1\ge 15$, which is false. So
$\langle x,y,z\rangle=W$: every subdesign of order $7$, $9$ or $13$ is the
closure of a triple of its own points. Closing each of the
$\binom{27}{3}=2925$ triples of points of $\mathcal D$ therefore produces every
such subdesign, and doing so yields exactly one closed set of size $7$, namely
$U$, and none of size $9$ or $13$. (The closed sets of size $3$ are the $117$
blocks.)

The set of points lying on a subdesign of order $7$ is preserved by every
automorphism, and so is the number of such subdesigns through a given point.
That number takes two values, $1$ and $0$, so $\operatorname{Aut}(\mathcal D)$
has more than one orbit on $V$.
\end{proof}

\section{A compact form}

The design is invariant under an automorphism of order $3$, which allows it to be
recorded by five of the thirteen classes. Let
\[
 \sigma=(1\ 2\ 4)(3\ 6\ 5)(9\ 10\ 11)(12\ 13\ 14)(15\ 16\ 17)
        (18\ 19\ 20)(21\ 22\ 23)(24\ 25\ 26),
\]
fixing $0,7,8$ and no other point. Then $\sigma$ has order $3$, it permutes the
$117$ triples of $\mathcal D$ among themselves, and on $U$ it acts as the
collineation $x\mapsto 2x \bmod 7$ of the Fano plane. Its fixed set
$\{0,7,8\}$ is itself a block; that block lies in $C_{01}$. On classes,
$\sigma$ induces
\[
 \pi=(C_{00}\ C_{06}\ C_{04})(C_{02}\ C_{03}\ C_{05})
     (C_{07}\ C_{08}\ C_{09})(C_{10}\ C_{11}\ C_{12}),
\]
fixing $C_{01}$. The five classes $C_{00},C_{01},C_{02},C_{07},C_{10}$ meet
every $\pi$-orbit, so applying $\sigma$ and $\sigma^2$ to them and placing the
images in the classes named by $\pi$ recovers all thirteen classes of
Table~\ref{tab:kts}. Five base classes therefore suffice to record the design;
Table~\ref{tab:kts} remains the definition.

\section{Scope}

A $\KTS(27)$ with an $\STS(7)$ subdesign exists; one is printed in
Table~\ref{tab:kts}. In the terms of \cite{DL} this is the cell $(v,u)=(27,7)$;
the orders $v=33$ and $v=39$ of the same sentence are untouched, and Open
Problem~4 of \cite{Stinson} is not settled here. No priority is claimed, and no
novelty either: whether a $\KTS(27)$ with an $\STS(7)$ subdesign has been
exhibited before, in \cite{KO}, in \cite[\S19.7]{ColbournRosa}, in the published
version of \cite{DL}, or in a design database, was not checked here, so this
note does not establish that any still-open request is answered by it. Nothing
is claimed about the number of such designs, about isomorphism (none was
rejected), or about the order of $\operatorname{Aut}(\mathcal D)$ beyond its
containing $\sigma$ and, by Proposition~\ref{prop:inv}, not being transitive on
points.

\section{Verification}

Both parts of Theorem~\ref{thm:main} are inspections of Table~\ref{tab:kts}, and
no input other than that table is used anywhere above.

A verification program \texttt{verify.py} accompanies this paper, together with
the transcript of a run. It parses Table~\ref{tab:kts} from the printed text,
in exact integer arithmetic and using only the Python standard library, and
re-derives every quantity stated here: the resolution and the $351$ pairs, the
seven inner blocks and their Fano structure, the $7+70+40$ census, the order and
action of $\sigma$ and the induced $\pi$, the reconstruction from the five base
classes, and the subdesign census of Proposition~\ref{prop:inv} by closing all
$2925$ point triples. All $59$ checks pass. The transcript also mentions a
second design found in the underlying search; it is not printed here and nothing
about it is claimed or verified.

\begin{thebibliography}{9}

\bibitem{CMM}
C.~J. Colbourn, S.~S. Magliveras, and R.~A. Mathon,
\emph{Transitive Steiner and Kirkman triple systems of order 27},
Math. Comp. \textbf{58} (1992), no.~197, 441--449.
\href{https://doi.org/10.1090/S0025-5718-1992-1106962-5}
{doi:10.1090/S0025-5718-1992-1106962-5}.

\bibitem{ColbournRosa}
C.~J. Colbourn and A.~Rosa,
\emph{Triple Systems},
Oxford Mathematical Monographs, Oxford University Press, 1999.

\bibitem{DL}
P.~Dukes and E.~R. Lamken,
\emph{An update on the existence of Kirkman triple systems with Steiner triple
systems as subdesigns},
J. Combin. Des. \textbf{30} (2022), no.~8, 581--608.
\href{https://doi.org/10.1002/jcd.21844}{doi:10.1002/jcd.21844}.
Quoted here from the e-print
\href{https://arxiv.org/abs/2110.07874v1}{arXiv:2110.07874v1}.

\bibitem{KO}
J.~I. Kokkala and P.~R.~J. \"Osterg\aa rd,
\emph{Kirkman triple systems with subsystems},
Discrete Math. \textbf{343} (2020), no.~9, Article 111960.
\href{https://doi.org/10.1016/j.disc.2020.111960}{doi:10.1016/j.disc.2020.111960}.

\bibitem{Stinson}
D.~R. Stinson,
\emph{A survey of Kirkman triple systems and related designs},
Discrete Math. \textbf{92} (1991), 371--393.

\end{thebibliography}

\end{document}
